\documentclass{rebuttal}
\renewcommand{\locallang}{PL}
\input{preamble.tex}


\title{Odpowiedź na pytania Recenzenta}
\author{Kacper Kania}
\date{}


\begin{document}
\maketitle

\responsedelimiter

\paragraph{Pytanie.}
  Z uwag polemicznych, metody proponowane w~rozprawie (a przynajmniej niektóre z
  nich, m.in. np. \clog, Sec. 5.3.3 i~5.3.4) wydają się wymagać dużego nakładu
  pracy badacza w~procesie uczenia (tzw. model babysitting).

\paragraph{Odpowiedź.}
  W naszych badaniach nie spotkaliśmy się z~problemem \textit{babysittingu}.
  Początkową trudnością okazało się zbudowanie wystarczająco ekspresywnych sieci
  MLP, które pozwalałyby modelować wartości Gaussów przy dowolnie zadanym
  poziomie szczegółów, a~jednocześnie pozwalały na uczenie przy wykorzystaniu
  pojedynczej karty A100 z~40GB pamięci VRAM. Pozostałe parametry, takie wagi
  regularyzacji, zostały dobrane empirycznie i nie podlegały znaczącym
  poszukiwaniom optymalnych wartości przy użyciu np. \textit{grid search}.
  Dodatkowo, raz ustalone wartości były wystarczające, aby osiągnąć
  satysfakcjonujące wyniki na dowolnym z~użytych w~publikacji zbiorów danych.

  \responsedelimiter

\paragraph{Pytanie.}
  Zastanawiam się także nad zdolnościami proponowanych metod do ekstrapolacji.
  Czy np. \conerf potrafi też ekstrapolować, raczej niż jedynie interpolować?
  (np. pomiędzy położeniami, orientacjami obiektów, ale także pomiędzy
  wartościami kontrolujących atrybutów). Termin ,,ekstrapolacja'' pojawiają się
  w~m.in. w~RQ2 i~Rozdziale~3, ale nie doczekał się chyba w~rozprawie
  definitywnych konkluzji.

\paragraph{Odpowiedź.}
  Zarówno \conerf oraz \blendfields zostały zaprojektowane z~myślą o~osiągnięciu
  jak najlepszej wizualnej jakości rekonstrukcji przy minimalnym nakładzie
  danych. W kontekście użytych metod możemy rozpatrywać ich zdolności do
  ekstrapolacji względnem pewnych osi. Dla przykładu, NeRFy mogą ekstrapolować
  jakość rekonstrukcji dla nowych widoków, które nie były widoczne w~czasie
  treningu (np. tył głowy). \conerf i~\blendfields są rozszerzeniami NeRFów,
  przez co naturalnie dziedziczą tę oś. Natomiast w~naszym przypadku, jesteśmy
  zainteresowani głównie ekstrapolacją w kontekście rzadkich sygnałów
  warunkujących naszych model. Dla CoNeRFa są to wartości atrybutów, natomiast
  BlendFieldsa---wartości wektora ekspresji. Zdolność do ekstrapolacji
  w~CoNeRFie osiągamy poprzez część modelu predykującą maskę 3D, której celem
  jest rozdzielenie wpływu wartości atrybutów na określone obszary w~przestrzeni
  NeRFa. Jednocześnie, ta część jest niezależna od samych wartości atrybutów
  i~działa w przestrzeni kanonicznej. Dzięki temu możemy dowolnie mieszać
  wartości atrybutów, mimo że w~czasie treningu przedstawione były tylko pewne
  możliwe kombinacje atrybutów. Na przykład, w~czasie treningu wektory atrybutów
  wynosiły $[1, 0, 0], [0, 1, 0], [0, 0, 1]$, a~dzięki naszemu systemu
  maskowania, możliwe jest realistyczne renderowanie obiektu zainteresowania
  przy podaniu wektora $[1, 1, 1]$ na wejście. Przykład takiej ekstrapolacji
  prezentuję na Rysunku~2 Rozdziału~2 oraz na wizualizacjach znajdujących się na
  stronie \url{https://conerf.github.io}. Natomiast w~\blendfields mamy do
  czynienia z~wielowymiarowym wektorem reprezentującym wartości ekspresji. Ten
  wektor warunkuje pozycje wierzchołków czworościanów, które budują naszą
  przestrzeń. Ze względu, że wartości kolorów predykowanych przez model są
  warunkowane tylko lokalnym sąsiedztwem czworościanów, punkty w~przestrzeni są
  niezależne od siebie. To powoduje, że nawet przy wykorzystaniu niewielu
  skrajnych wyrazów twarzy, jesteśmy w~stanie z~wysoką jakością przedstawiać
  wygląd zmarszczek nawet przy nowych wyrazach twarzy. Taką ekstrapolacją jest
  właśnie na renderze po prawej stronie Rysunku~17 Rozdziału~3. Dodatkowo,
  prezentuję nie widziane w~czasie treningu wyrazy twarzy i~ich wpływ na
  predykowany kolor skóry na stronie \url{https://blendfields.github.io}.

  \responsedelimiter

\end{document}